\section{L'identità segreta degli oggetti: Prototipi e Classi}

\subsection{Recap: Valori, Riferimenti e Indirezione}
In JavaScript, la distinzione tra tipi base e complessi è fondamentale:
\begin{itemize}
    \item \textbf{Tipi Base:} (booleani, numeri, stringhe) $\rightarrow$ Semantica per \textbf{valore}.
    \item \textbf{Tipi Complessi:} (oggetti, array, funzioni) $\rightarrow$ Semantica per \textbf{riferimento}.
\end{itemize}

La semantica per riferimento utilizza indirizzi di memoria come valori. L'accesso ai dati avviene tramite operazioni di \textbf{indirezione} (spesso implicite in JS).

\begin{definition}[Operazioni di Indirezione]
    \begin{itemize}
        \item \texttt{oggetto.chiave}: Va in memoria all'indirizzo dell'oggetto, cerca la chiave e restituisce il contenuto.
        \item \texttt{array[indice]}: Analogo al precedente, usando l'indice come chiave.
        \item \texttt{funzione(...)}: Va in memoria a recuperare la definizione ed esegue il codice.
    \end{itemize}
\end{definition}

\section{Identità vs Uguaglianza: Il caso dei Grafi}
L'uso dei riferimenti è cruciale nelle strutture dati complesse come i grafi, dove lo \textbf{stesso} oggetto nodo può essere puntato da molteplici archi.

\begin{concept}[Identità]
Dire "lo stesso oggetto" (Identità) non è come dire "un oggetto uguale" (Uguaglianza).
\begin{itemize}
    \item \textbf{Uguaglianza:} Due oggetti distinti con lo stesso contenuto (es. \texttt{\{val:1\}} e \texttt{\{val:1\}}).
    \item \textbf{Identità:} Due riferimenti che puntano alla stessa area di memoria.
\end{itemize}
\end{concept}

Per modellare un grafo $G=\langle N,E\rangle$ dove più archi insistono sullo stesso nodo, è indispensabile usare variabili che contengano i riferimenti ai nodi, e riutilizzare quelle variabili nella creazione degli archi.

\section{Costruttori e l'Operatore \texttt{new}}
Per garantire che gli oggetti abbiano sempre la struttura corretta (campi giusti), si usano funzioni \textbf{costruttori}.

\subsection{Funzionamento di \texttt{new}}
L'operatore \texttt{new}, anteposto a una chiamata di funzione, altera il comportamento standard:
\begin{enumerate}
    \item \textbf{Allocazione:} Crea un nuovo oggetto vuoto (chiamiamolo \textit{o}).
    \item \textbf{Invocazione:} Chiama la funzione costruttore impostando \texttt{this = o}.
    \item \textbf{Arricchimento:} Il codice della funzione usa \texttt{this} per aggiungere proprietà e metodi a \textit{o}.
    \item \textbf{Ritorno:} Restituisce implicitamente \textit{o}.
\end{enumerate}

\begin{jsbox}[Confronto Costruttori]
// 1. Factory Function (Sconsigliato - crea oggetti, non istanze legate)
function grafo(N, E) {
    return { nodi: N, archi: E, size() { return this.nodi.length } };
}

// 2. Constructor Function con 'new' (Standard pre-ES6)
function Grafo(N, E) {
    this.nodi = N;
    this.archi = E;
    // this è l'oggetto appena creato da new
}
var g = new Grafo([], []);
\end{jsbox}

\section{I Prototipi}
Il meccanismo dei prototipi permette agli oggetti di ereditare funzionalità senza duplicarle in memoria.

\begin{definition}[Algoritmo di Lookup delle Proprietà]
Quando si accede a \texttt{obj.key}:
\begin{enumerate}
    \item JS cerca \texttt{key} direttamente in \texttt{obj}. Se trovata, restituisce il valore.
    \item Se non trovata, cerca nel \textbf{prototipo} di \texttt{obj}.
    \item Continua a risalire la catena dei prototipi.
    \item Se arriva a un prototipo \texttt{null} (fine catena) e non ha trovato nulla, restituisce \texttt{undefined}.
\end{enumerate}
\textit{Nota: La scrittura avviene sempre sull'oggetto stesso, mai sul prototipo.}
\end{definition}

\subsection{Visualizzazione della Catena}
Possiamo ispezionare il prototipo con \texttt{Object.getPrototypeOf(o)} o la proprietà deprecata \texttt{\_\_proto\_\_}.

\begin{center}
\begin{tikzpicture}[
    object/.style={rectangle, draw=paletteIndigo, fill=paletteLavender!20, thick, minimum width=2.5cm, minimum height=1.5cm, align=center},
    pointer/.style={->, >=stealth, thick, color=paletteRegalia},
    node distance=1.5cm and 2cm
]

% Oggetto array
\node[object] (arr) {Array \texttt{a} \\ \texttt{[1, 5, 7]} \\ \footnotesize \texttt{\_\_proto\_\_: ?}};

% Array Prototype
\node[object, right=of arr] (arrProto) {\textbf{Array.prototype} \\ \textit{push, pop, map...} \\ \footnotesize \texttt{\_\_proto\_\_: ?}};

% Object Prototype
\node[object, right=of arrProto] (objProto) {\textbf{Object.prototype} \\ \textit{toString, valueOf...} \\ \footnotesize \texttt{\_\_proto\_\_: null}};

% Null
\node[draw=none, right=of objProto] (null) {\textbf{null}};

% Collegamenti
\draw[pointer] (arr.east) -- node[above]{\footnotesize prototipo} (arrProto.west);
\draw[pointer] (arrProto.east) -- node[above]{\footnotesize prototipo} (objProto.west);
\draw[pointer] (objProto.east) -- (null.west);

\end{tikzpicture}
\end{center}

\begin{warning}[Enumerabilità]
Il ciclo \texttt{for (k in o)} e \texttt{Object.keys(o)} restituiscono solo le chiavi \textbf{proprie} dell'oggetto, ignorando quelle ereditate dal prototipo.
\end{warning}

\section{Prototipi e Costruttori: Il Doppio Salto Carpiato}
Esiste un legame simbiotico tra una funzione costruttore e i prototipi degli oggetti che essa crea.
\begin{concept}[Proprietà \texttt{.prototype}]
Ogni funzione \texttt{F} ha una proprietà \texttt{F.prototype}. Quando chiamiamo \texttt{new F()}, l'oggetto creato riceve automaticamente \texttt{F.prototype} come suo prototipo interno (\texttt{\_\_proto\_\_}).
\end{concept}

Questo permette di definire metodi una sola volta nel prototipo del costruttore, rendendoli disponibili a tutte le istanze:
\begin{jsbox}[Pattern Costruttore + Prototipo]
function Persona(n, e) {
    this.nome = n; // Stato interno (diverso per ogni istanza)
    this.età = e;
}
// Comportamento condiviso (uguale per tutte le istanze)
Persona.prototype.compleanno = function() { this.età++; };

var p = new Persona("Pippo", 35);
// p.__proto__ === Persona.prototype è TRUE
\end{jsbox}

\section{Le Classi (Syntactic Sugar)}
Le versioni recenti di JavaScript introducono la keyword \texttt{class}. Non è un nuovo modello a oggetti, ma una sintassi più pulita sopra il sistema prototipale esistente.

\begin{jsbox}[Sintassi Class]
class Persona {
    constructor(nome, eta) {
        this.nome = nome;
        this.età = eta;
    }

    compleanno() {
        this.età++;
    }
}
\end{jsbox}

\subsection{Cosa accade dietro le quinte?}
Anche usando \texttt{class}:
\begin{itemize}
    \item \texttt{typeof Persona} restituisce \textbf{"function"}.
    \item I metodi definiti nella classe (es. \texttt{compleanno}) finiscono in \texttt{Persona.prototype}.
    \item Le proprietà definite nel costruttore (es. \texttt{nome}) finiscono direttamente nell'istanza oggetto.
\end{itemize}

\section{La "Mappa Proibita" (Struttura Profonda di JS)}
Le slide finali mostrano la struttura ricorsiva e intrecciata di JS. Ecco la logica:

\begin{itemize}
    \item \textbf{Ogni funzione è un oggetto:} Quindi \texttt{Array} (il costruttore) ha come prototipo \texttt{Function.prototype}.
    \item \textbf{Function.prototype è un oggetto:} Quindi il suo prototipo è \texttt{Object.prototype}.
    \item \textbf{Object è una funzione:} Il suo prototipo è \texttt{Function.prototype}.
\end{itemize}

\begin{center}
\begin{tikzpicture}[
    node distance=2cm,
    entity/.style={rectangle, draw=paletteRegalia, fill=white, thick, minimum size=1cm},
    proto/.style={->, >=stealth, dashed, color=gray!80}
]

\node[entity] (funcProto) {Function.prototype};
\node[entity, below=of funcProto] (objProto) {Object.prototype};
\node[entity, right=of funcProto] (ArrayConst) {Array (Funzione)};
\node[entity, below=of ArrayConst] (ArrProto) {Array.prototype};
\node (null) [right=of objProto] {null};

% Relazioni
\draw[proto] (ArrayConst) -- node[right]{\tiny \_\_proto\_\_} (funcProto);
\draw[proto] (funcProto) -- node[right]{\tiny \_\_proto\_\_} (objProto);
\draw[proto] (ArrProto) -- node[right]{\tiny \_\_proto\_\_} (objProto);
\draw[proto] (objProto) -- (null);

% Relazione prototype
\draw[->, thick, color=paletteIndigo] (ArrayConst) -- node[left]{\tiny .prototype} (ArrProto);

\end{tikzpicture}
\end{center}
